\section{Detectors}
\label{sec:models}

We compare four detector families under shared splits, evaluator and bootstrap. Table~\ref{tab:recipes} lists the training settings.

\subsection{YOLO12-L}
The stock Ultralytics YOLO12-L~\cite{ultralytics,yolov12} is COCO-pretrained and uses the CIoU loss~\cite{diou} and AdamW~\cite{adamw} (learning rate 0.001, weight decay 0.01) with a cosine schedule~\cite{sgdr}. Training uses mosaic~\cite{yolov4}, vertical and horizontal flips and random quarter turns. A quarter turn maps an axis-aligned box onto an exactly tight box, and drone views have no canonical ``up''.

\subsection{BladeYOLO-v2}
\label{sec:byv2}
BladeYOLO-v2 is our re-engineering of the BladeYOLO idea after a code audit (Section~\ref{sec:history}). It keeps the COCO-pretrained YOLOv12-L backbone, neck and head, so that 1{,}239 pretrained tensors transfer, and adds BladeYOLO's ingredients as zero-initialised residual branches (Fig.~\ref{fig:arch-byv2}). Training therefore starts exactly at the pretrained detector, and it learns how much of each addition to use:
\begin{itemize}
\item \textbf{DinoFuse.} Frozen DINOv3 features~\cite{dinov3} from blocks 4, 8 and 12 (ViT-S/16 by default; ViT-B/16 and ViT-L/16 in the ablations) are resized to the P3, P4 and P5 backbone levels. They are projected there and added through a zero-initialised $1\times1$ convolution: $X_\ell' = X_\ell + W_\ell\,\phi(\mathrm{Resize}(D_{b(\ell)}))$ with $W_\ell=0$ at initialisation.
\item \textbf{Deformable refinement} of the P3 and P4 neck outputs with modulated deformable convolution (DCNv2~\cite{dcnv2}), for irregular contours such as cracks and erosion edges. Offsets start at zero and modulation at 0.5.
\item \textbf{Selective-scan refinement} of the P5 neck output: a four-direction VMamba-style scan~\cite{vmamba,mamba} in the chunked Mamba-2 form, written in plain PyTorch. It matches a step-by-step reference recurrence to $1.8\times10^{-14}$.
\item \textbf{WIoU v3}~\cite{wiou} box loss in place of CIoU. The assigner keeps CIoU.
\end{itemize}
With ViT-S/16 the model has 53.7\,M parameters (32.1\,M trainable) and 233.5 GFLOPs. Deformable and selective-scan blocks run in fp32 under mixed precision~\cite{mixedprecision}, and YOLOv12's area-attention scores are computed in fp32. Without these changes, fp16 overflows produced NaNs at epochs 9, 17 and 99--149 (Appendix~\ref{app:engineering}).

\begin{figure*}[!t]
\centering
\resizebox{0.92\textwidth}{!}{% BladeYOLO-v2 (models/bladeyolo2-l.yaml; models/dino_fusion.py, morphology.py, ss2d.py, wiou.py)
\begin{tikzpicture}[
  font=\scriptsize, >={Stealth[length=4pt]}, node distance=3mm and 4mm,
  blk/.style={draw=black!70, rounded corners=1.5pt, align=center, minimum height=7mm, inner sep=2pt, fill=#1},
  blk/.default=white,
  lab/.style={font=\tiny, text=black!70, align=center},
  frozen/.style={blk=cyan!10, dashed, thick},
  coco/.style={blk=orange!12},
  new/.style={blk=green!12},
]
% backbone row
\node[blk=gray!8, minimum width=10mm] (img) {Image\\640 px};
\node[coco, right=of img, minimum width=13mm] (s1) {Stem\\P1, P2};
\node[coco, right=of s1, minimum width=11mm] (p3) {C3k2\\P3 (1/8)};
\node[coco, right=of p3, minimum width=11mm] (p4) {A2C2f\\P4 (1/16)};
\node[coco, right=of p4, minimum width=11mm] (p5) {A2C2f\\P5 (1/32)};
\draw[->] (img) -- (s1); \draw[->] (s1) -- (p3); \draw[->] (p3) -- (p4); \draw[->] (p4) -- (p5);
% DINOv3 and fusion
\foreach \k in {3,4,5} {
  \node[new, below=4mm of p\k, minimum width=11mm, minimum height=5.5mm] (f\k) {DinoFuse};
  \draw[->] (p\k) -- (f\k);
}
\node[frozen, below=9mm of f4, minimum height=6mm] (dino) {DINOv3 ViT-S/16 or ViT-B/16, \textbf{frozen}};
\draw[->] (img) |- (dino.west);
\foreach \k/\b in {3/4, 4/8, 5/12} {
  \draw[->, cyan!50!black] (dino.north -| f\k) -- node[right, lab, pos=0.25]{blk \b} (f\k.south);
}
% neck
\node[coco, right=9mm of p5, minimum width=19mm, minimum height=27mm, yshift=-8mm] (neck) {YOLOv12-L neck\\(PAN, A2C2f)\\top-down and\\bottom-up};
\foreach \k/\d in {5/2.2, 4/4.2, 3/6.2} {
  \draw[->, thick] (f\k.east) -- ++(1.4mm,0) |- ([yshift=-\d mm]f5.south -| neck.west);
}
% refinements after the neck
\node[new, right=8mm of neck, minimum width=13mm, yshift=6mm] (r5) {SS2D\\4-dir scan};
\node[new, below=1.5mm of r5, minimum width=13mm] (r4) {DCNv2\\refine};
\node[new, below=1.5mm of r4, minimum width=13mm] (r3) {DCNv2\\refine};
\draw[->] ([yshift=6mm]neck.east) -- node[above, lab]{P5} (r5.west);
\draw[->] ([yshift=-2.5mm]neck.east) -- node[above, lab]{P4} (r4.west);
\draw[->] ([yshift=-11mm]neck.east) -- node[above, lab]{P3} (r3.west);
\node[coco, right=6mm of r4, minimum width=13mm, minimum height=24mm] (head) {Detect\\head\\(DFL,\\classes)};
\foreach \k in {3,4,5} {\draw[->] (r\k.east) -- (r\k.east -| head.west);}
\node[blk=gray!8, below=3mm of head, minimum width=13mm] (loss) {WIoU v3\\box loss};
\draw[->, dashed] (head) -- (loss);
% legend
\node[frozen, minimum height=3mm, minimum width=5mm, below=3mm of dino.south west, anchor=north west] (l1) {};
\node[right=0.5mm of l1, lab, anchor=west] {frozen};
\node[coco, minimum height=3mm, minimum width=5mm, right=9mm of l1] (l2) {};
\node[right=0.5mm of l2, lab, anchor=west] {COCO-pretrained YOLOv12-L};
\node[new, minimum height=3mm, minimum width=5mm, right=33mm of l2] (l3) {};
\node[right=0.5mm of l3, lab, anchor=west] {zero-initialised residual additions};
\end{tikzpicture}
}
\caption{BladeYOLO-v2. A COCO-pretrained YOLOv12-L with three kinds of zero-initialised residual additions: frozen DINOv3 features fused into the backbone levels P3--P5 (DinoFuse), DCNv2 refinement of the P3/P4 neck outputs, and a four-direction selective scan (SS2D) on the P5 neck output. The box loss is WIoU v3. At initialisation the network computes exactly the pretrained YOLOv12-L. Section~\ref{sec:ablation} shows that the frozen DINOv3 features carry its entire advantage.}
\label{fig:arch-byv2}
\end{figure*}

\subsection{RF-DETR}
RF-DETR-M~\cite{rfdetr} (\texttt{rfdetr} 1.11.2) builds on LW-DETR~\cite{lwdetr}. It uses a windowed DINOv2-S/16 encoder~\cite{dinov2}, fine-tuned end to end, whose blocks 3, 6, 9 and 12 feed a single-level (P4, stride 16) projector with 256 channels. A two-stage head selects 300 queries from this memory, and four deformable cross-attention decoder layers~\cite{deformabledetr} refine them. Training uses group-DETR one-to-many assignment (13 groups) and IoU-aware BCE classification with GIoU and L1 box losses~\cite{giou}. We fine-tune from the COCO checkpoint with AdamW (base learning rate $10^{-4}$, encoder $1.5\times10^{-4}$ with layer-wise decay 0.8, weight decay $10^{-4}$), gradient clipping at 0.1 and per-batch multi-scale training. The effective batch is 16 (4 $\times$ 4 accumulation). An exponential moving average~\cite{polyak} of the weights (decay 0.993) is evaluated every epoch, and the best EMA checkpoint on validation is kept. RF-DETR's default learning-rate step at epoch 100 never fires in a 100-epoch run, so we use an explicit cosine schedule (``G26 recipe'').

\subsection{RF-DETR-FP: A Frozen Foundation Prior Beside a Fine-Tuned Encoder}
\label{sec:fp}

\textbf{Motivation.} Three findings of this study, each made before RF-DETR-FP was designed, point to the same lever:
\begin{enumerate}
\item Frozen DINOv3 is the one component that ever helped in our YOLO-family ablations. It adds $+3.9$, $+8.8$ and $+9.0$ \mcoco at 25, 50 and 100\% labels, and ViT-B beats ViT-S (Section~\ref{sec:ablation}).
\item RF-DETR lacks what DINOv3 carries. On WTBlade-group, the frozen-DINOv3 ViT-B YOLO beats fine-tuned RF-DETR-M (18.0 against 16.2/16.9). On WTBD, DINOv3 ViT-B box embeddings re-rank RF-DETR's own outputs for a gain (H18, Table~\ref{tab:ledger}), so its ranking deficit is information that DINOv3 holds and RF-DETR does not use.
\item That information cannot simply be written into RF-DETR's encoder. Fine-tuning moves RF-DETR's features away from generic neighbourhoods (G22-O). Further self-supervised adaptation of its encoder halved the effective rank of its last stage (G27 D1, Section~\ref{sec:failed}), and multi-dataset pretraining did not help (D2).
\end{enumerate}
We therefore keep a second, \emph{frozen} encoder beside the fine-tuned one, rather than adapting either (Fig.~\ref{fig:arch-fp}).

\begin{figure*}[!t]
\centering
\resizebox{\textwidth}{!}{% RF-DETR-M with a frozen DINOv3 prior (G29, "FP"). tools/g29_prior.py
\begin{tikzpicture}[
  font=\scriptsize, >={Stealth[length=4pt]}, node distance=4mm and 5mm,
  blk/.style={draw=black!70, rounded corners=1.5pt, align=center, minimum height=8mm, inner sep=2pt, fill=#1},
  blk/.default=white,
  op/.style={draw=black!70, circle, inner sep=0.5pt, minimum size=4.2mm, fill=white},
  lab/.style={font=\tiny, text=black!70, align=center},
  frozen/.style={blk=cyan!10, dashed, thick},
  train/.style={blk=orange!12},
  new/.style={blk=green!12},
]
% ---------- main RF-DETR path ----------
\node[blk=gray!8, minimum width=12mm] (img) {Image\\[3pt]$I\in\mathbb{R}^{3\times H\times W}$};
\node[train, right=6mm of img, minimum width=22mm] (enc) {RF-DETR encoder\\windowed DINOv2-S/16\\(fine-tuned)};
\node[train, right=of enc, minimum width=15mm] (proj) {Projector\\$C{=}256$, stride 16};
\node[op, right=7mm of proj] (add) {$+$};
\node[train, right=7mm of add, minimum width=20mm] (qs) {Two-stage\\query selection\\(top-$K$ proposals)};
\node[train, right=of qs, minimum width=20mm] (dec) {Deformable\\cross-attention\\decoder};
\node[blk=gray!8, right=of dec, minimum width=12mm] (out) {Boxes,\\classes};
\draw[->] (img) -- (enc);
\draw[->] (enc) -- (proj);
\draw[->] (proj) -- node[above, lab]{$F$} (add);
\draw[->] (add) -- node[above, lab]{$F'$} (qs);
\draw[->] (qs) -- (dec);
\draw[->] (dec) -- (out);
\draw[->, black!60] ($(add.east)!0.5!(qs.west)$) |- ([yshift=-3.5mm]qs.south) -| node[pos=0.75, right, lab]{memory $F'$} (dec.south);

% ---------- frozen prior path ----------
\node[frozen, below=17mm of enc, minimum width=22mm] (dino) {DINOv3 ViT-B/16\\\textbf{frozen}, no grad};
\node[blk=cyan!5, right=of dino, minimum width=15mm] (cat) {blocks 4, 8, 12\\LN, concat\\$P\in\mathbb{R}^{2304\times\frac{H}{16}\times\frac{W}{16}}$};
\draw[->] (img) |- (dino);
\draw[->] (dino) -- (cat);

% ---------- PriorFuse ----------
\node[new, right=8mm of cat, minimum width=9mm] (red) {LN\\$1{\times}1$\\Reduce\\GELU};
\node[op, right=3mm of red] (cc) {\tiny $\Vert$};
\node[new, right=3mm of cc, minimum width=9mm] (mix) {$1{\times}1$\\Mix\\LN\\GELU};
\node[new, right=3mm of mix, minimum width=9mm] (dw) {DW\\$3{\times}3$\\GELU};
\node[new, right=3mm of dw, minimum width=11mm] (o) {$1{\times}1$ Out\\$W_o{=}0$\\$b_o{=}0$};
\draw[->] (cat) -- (red);
\draw[->] (red) -- (cc);
\draw[->] (cc) -- (mix);
\draw[->] (mix) -- (dw);
\draw[->] (dw) -- (o);
\draw[->] (proj.south) ++(2mm,0) |- ([yshift=4mm]cc.north) -- (cc.north);
\draw[->, green!40!black, thick] (o.north) |- ([yshift=-7mm]add.south) -- node[left, lab]{$\Delta F$} (add.south);
\begin{scope}[on background layer]
  \node[draw=green!50!black, dotted, thick, rounded corners=3pt, fit=(red)(cc)(mix)(dw)(o), inner sep=2.2mm,
        label={[lab, anchor=north]south:PriorFuse (0.79\,M trainable parameters; identity at initialisation)}] (pf) {};
\end{scope}
% legend
\node[frozen, minimum height=3mm, minimum width=5mm, below=5mm of dino.south west, anchor=north west] (l1) {};
\node[right=0.5mm of l1, lab, anchor=west] {frozen, outside optimizer, EMA and checkpoint};
\node[train, minimum height=3mm, minimum width=5mm, below=1.5mm of l1] (l2) {};
\node[right=0.5mm of l2, lab, anchor=west] {RF-DETR-M, fine-tuned};
\node[new, minimum height=3mm, minimum width=5mm, below=1.5mm of l2] (l3) {};
\node[right=0.5mm of l3, lab, anchor=west] {new layers};
\end{tikzpicture}
}
\caption{RF-DETR-FP. A frozen DINOv3 ViT-B/16 runs beside RF-DETR-M's fine-tuned windowed DINOv2-S encoder on the same normalised input. Both have patch 16, so the frozen features from blocks 4, 8 and 12 (2{,}304 channels) lie on the projector's grid. PriorFuse adds them to the projector output $F$ before two-stage query selection and decoding. Its output convolution is zero-initialised, so the model at step 0 is exactly RF-DETR-M. DINOv3 sits outside the optimiser, the EMA copy and the checkpoint. Only PriorFuse (0.79\,M parameters) is new.}
\label{fig:arch-fp}
\end{figure*}

\textbf{Architecture.} Let $I\in\mathbb{R}^{3\times H\times W}$ be the ImageNet-normalised input, which is also DINOv3's normalisation. Let $F=\mathrm{Proj}(E_\theta(I))\in\mathbb{R}^{C\times\frac{H}{16}\times\frac{W}{16}}$, with $C=256$, be RF-DETR's projector output. The frozen prior is
\begin{equation}
P=\big[\,\mathrm{LN}(\Phi_4(I));\ \mathrm{LN}(\Phi_8(I));\ \mathrm{LN}(\Phi_{12}(I))\,\big]\in\mathbb{R}^{2304\times\frac{H}{16}\times\frac{W}{16}},
\label{eq:prior}
\end{equation}
where $\Phi_k$ is the patch-token map after block $k$ of DINOv3 ViT-B/16, computed without gradient, and LN is DINOv3's final layer norm~\cite{layernorm}. PriorFuse computes
\begin{align}
S &= \mathrm{GELU}\big(W_r\ast\mathrm{LN}(P)\big), \label{eq:reduce}\\
Z &= \mathrm{GELU}\big(\mathrm{LN}(W_m\ast[F;S])\big), \label{eq:mix}\\
F' &= F + W_o\ast\mathrm{GELU}\big(\mathrm{DW}_{3\times3}(Z)\big),\qquad W_o=0,\ b_o=0\ \text{at}\ t=0, \label{eq:fuse}
\end{align}
where $W_r$ ($2304\rightarrow C$), $W_m$ ($2C\rightarrow C$) and $W_o$ ($C\rightarrow C$) are $1\times1$ convolutions, $\mathrm{DW}_{3\times3}$ is a depthwise convolution~\cite{xception} and GELU is from~\cite{gelu}. $F'$ replaces $F$ as the memory from which the two-stage head selects queries and which the decoder attends to. Because the grids coincide on RF-DETR-M, no resampling is needed. For RF-DETR sizes with other projector levels, $P$ is resized bilinearly.

\textbf{Properties.} (i) At initialisation the patched backbone returns bit-identical features to RF-DETR, which we verify in a unit test at 320 and 768\,px. (ii) The frozen model lives in a process-wide cache, not as a submodule, so it is not in \texttt{parameters()}, the optimiser, the EMA copy or \texttt{state\_dict()}. Checkpoints carry only the 0.79\,M fusion parameters, and gradients reach PriorFuse and never DINOv3, which a unit test also verifies. (iii) The recipe is otherwise unchanged: each FP run uses exactly the arguments of a baseline twin with the same hardware type, so that every comparison is paired. (iv) On a Tesla T4 the frozen ViT-B forward, plus a second job sharing the machine, costs 1.5--1.7$\times$ the baseline's wall time (Table~\ref{tab:recipes}).

\textbf{Novelty, stated plainly.} Fusing frozen foundation features into a detector is prior art: BladeYOLO and DINO-YOLO do it in YOLO~\cite{bladeyolo,dinoyolo}, and DEIMv2 puts DINOv3 in a DETR~\cite{deimv2}. What is specific here is narrower. First, a frozen generic encoder is placed \emph{beside} a fine-tuned one, chosen because of measured evidence that the fine-tuned one drifts from it and collapses when pushed further. Second, the test is leak-aware: we ask whether the gain is still there on images with no training twin. This is a defensible design, not a new module, and we evaluate it as such.

\begin{table*}[!t]
\centering
\caption{Training settings. All RF-DETR runs fine-tune the COCO checkpoint, select the best EMA checkpoint on validation and score test once. Wall times are per run on one Tesla T4 (fp16) unless noted.}
\label{tab:recipes}
\footnotesize
\begin{tabularx}{\textwidth}{@{}l l l l l X@{}}
\toprule
Model & Data & Input & Epochs & Batch & Notes \\
\midrule
YOLO12-L, BladeYOLO-v2 & WTBD-group / WTBD-official & 640 & 150 / 120 & 16 & AdamW 0.001, wd 0.01, cosine, mosaic, flips, Rot90; Ultralytics \\
YOLO12-L, BladeYOLO-v2 ViT-B & WTBlade-group & 640 & 50 & 16 & as above \\
RF-DETR-M (A0) & WTBD-group & 768 & 100 & 4$\times$4 & cosine; 2.5\,h (FP: 4.2\,h) \\
RF-DETR-M (A0) & WTBlade-group & 640 & 60 & 4$\times$4 & cosine; $\approx$3.5\,h (FP: 5.5\,h) \\
RF-DETR-M (G26) & WTBD-official & 768 & 100 & 2$\times$8 & cosine; 3\,h (FP: 4.9\,h) \\
RF-DETR-M (G26) & Wind Surface (2{,}570/286/714) & 640 & 100 & 4$\times$4 & cosine; 7\,h (FP: 10.8\,h) \\
RF-DETR-M (G26) & WTBlade-Defect official & 640 & 40 & 4$\times$4 & cosine; A100 (bf16) \\
RF-DETR-M (V1) & 25\% of WTBD-group, 4 subsets & 576 & 100 & 4$\times$4 & RF-DETR's constant schedule; 35\,min (FP: $\approx$1\,h) \\
RF-DETR-M class-agnostic (G28) & ID / X pools & 640 & 100, 60 / 12 & 4$\times$4 & cosine; X selected on its own held-out 5\% \\
\bottomrule
\end{tabularx}
\end{table*}

\textbf{Hardware.} Runs used Kaggle Tesla T4 GPUs (fp16 mixed precision) and one A100 (bf16). Every RF-DETR-FP run and its twin ran on T4s. G27's seed-1 runs ran on the A100 and its seed-0 runs on Kaggle, so every pair of arms within a seed shares hardware. The G29 queue ran two lanes in parallel per Kaggle session, one per T4, under a time guard. The guard starts no run with less than 30\,min left of an 11.3\,h budget, stops training 15\,min before the budget ends and never predicts from a stopped run. All ten G29 runs finished inside the guard.
